% GENERATED FILE. Edit canonical lessons, then run npm run generate:books.
% canonical-source-hash: fnv1a64:b43f793b52bee1a2
% canonical-lessons: IT-C452-tabaccheria, IT-C452-tabaccaio, IT-C452-edicola, IT-C452-edicolante, IT-C452-cartoleria, IT-R452-ripasso-445-452

\chapter{Giornali e tabacchi}
\label{ch:it-giornali-e-tabacchi}
\hlchaptermodality{452}

\begin{chapteropening}
\textbf{By the end of this chapter:} \emph{I can use small everyday shops: the tobacconist and his shop, the newsstand, the newsagent and the stationer.}

Complete the last lesson of chapter 452: I can use small everyday shops: the tobacconist and his shop, the newsstand, the newsagent and the stationer.
\end{chapteropening}

\section[la tabaccheria]{la tabaccheria — the tobacconist's}

\label{lesson:IT-C452-tabaccheria}



\begin{quote}
Before the new one: say the Italian for to peek, then the Italian for the shop window.
\end{quote}

\subsection*{You'll want to know: la tabaccheria}
\textbf{la tabaccheria} — \textquotedblleft{}the tobacconist's\textquotedblright{}.

\begin{quote}
In tabaccheria compro i biglietti dell'autobus.
\end{quote}

\begin{quote}
\emph{At the tobacconist's I buy bus tickets.}
\end{quote}

Marked by a big T sign; it sells tickets, stamps and phone credit too.

\begin{cousinweb}
\emph{tabaccheria} is built on \emph{tabacco} 'tobacco', which was borrowed from Spanish \emph{tabaco}, a word of Caribbean origin.
\end{cousinweb}

\begin{grammarlens}[title={What you've built}]
One more word for telling, explaining and describing at B1.
\end{grammarlens}

\subsection*{Your turn}
\begin{itemize}
\raggedright
  \item \emph{Say it:} \emph{la tabaccheria}
  \item \emph{Say it:} \emph{la tabaccheria}, once more
  \item \emph{Say it:} a sentence of your own with \emph{la tabaccheria}
  \item \emph{Recall:} say the Italian for to peek, then the Italian for the shop window, then say \emph{la tabaccheria} again
\end{itemize}

\begin{tcolorbox}[breakable,colback=teal!4,colframe=teal!35!black,title={Before you move on}]
What does la tabaccheria mean? (\textquotedblleft{}The tobacconist's\textquotedblright{}.) Say it once more.
\end{tcolorbox}
\section[il tabaccaio]{il tabaccaio — the tobacconist}

\label{lesson:IT-C452-tabaccaio}



\begin{quote}
Before the new one: say the Italian for the shop window, then the Italian for the tobacconist's.
\end{quote}

\subsection*{You'll want to know: il tabaccaio}
\textbf{il tabaccaio} — \textquotedblleft{}the tobacconist\textquotedblright{}.

\begin{quote}
Il tabaccaio mi dà il resto.
\end{quote}

\begin{quote}
\emph{The tobacconist gives me my change.}
\end{quote}

The person; the shop is \emph{la tabaccheria}.

\begin{cousinweb}
\emph{tabaccaio} is built on \emph{tabacco} 'tobacco', which was borrowed from Spanish \emph{tabaco}, a word of Caribbean origin.
\end{cousinweb}

\begin{grammarlens}[title={What you've built}]
One more word for telling, explaining and describing at B1.
\end{grammarlens}

\subsection*{Your turn}
\begin{itemize}
\raggedright
  \item \emph{Say it:} \emph{il tabaccaio}
  \item \emph{Say it:} \emph{il tabaccaio}, once more
  \item \emph{Say it:} a sentence of your own with \emph{il tabaccaio}
  \item \emph{Recall:} say the Italian for the shop window, then the Italian for the tobacconist's, then say \emph{il tabaccaio} again
\end{itemize}

\begin{tcolorbox}[breakable,colback=teal!4,colframe=teal!35!black,title={Before you move on}]
What does il tabaccaio mean? (\textquotedblleft{}The tobacconist\textquotedblright{}.) Say it once more.
\end{tcolorbox}
\section[l'edicola]{l'edicola — the newsstand, the kiosk}

\label{lesson:IT-C452-edicola}



\begin{quote}
Before the new one: say the Italian for the tobacconist's, then the Italian for the tobacconist.
\end{quote}

\subsection*{You'll want to know: l'edicola}
\textbf{l'edicola} — \textquotedblleft{}the newsstand, the kiosk\textquotedblright{}.

\begin{quote}
All'edicola compro il giornale ogni mattina.
\end{quote}

\begin{quote}
\emph{At the newsstand I buy the paper every morning.}
\end{quote}

It also means a small shrine with a statue in a wall.

\begin{cousinweb}
\emph{edicola} is a learned borrowing from Latin \emph{aedicula} 'small shrine' (EtymDB), from \emph{aedēs} 'temple, dwelling': the kiosk looks like a little shrine.
\end{cousinweb}

\begin{grammarlens}[title={What you've built}]
One more word for telling, explaining and describing at B1.
\end{grammarlens}

\subsection*{Your turn}
\begin{itemize}
\raggedright
  \item \emph{Say it:} \emph{l'edicola}
  \item \emph{Say it:} \emph{l'edicola}, once more
  \item \emph{Say it:} a sentence of your own with \emph{l'edicola}
  \item \emph{Recall:} say the Italian for the tobacconist's, then the Italian for the tobacconist, then say \emph{l'edicola} again
\end{itemize}

\begin{tcolorbox}[breakable,colback=teal!4,colframe=teal!35!black,title={Before you move on}]
What does l'edicola mean? (\textquotedblleft{}The newsstand, the kiosk\textquotedblright{}.) Say it once more.
\end{tcolorbox}
\section[l'edicolante]{l'edicolante — the newsagent}

\label{lesson:IT-C452-edicolante}



\begin{quote}
Before the new one: say the Italian for the tobacconist, then the Italian for the newsstand.
\end{quote}

\subsection*{You'll want to know: l'edicolante}
\textbf{l'edicolante} — \textquotedblleft{}the newsagent\textquotedblright{}.

\begin{quote}
L'edicolante mi tiene da parte la rivista.
\end{quote}

\begin{quote}
\emph{The newsagent keeps the magazine aside for me.}
\end{quote}

The person who runs \emph{l'edicola}.

\begin{cousinweb}
\emph{edicolante} is built on \emph{edicola}, a learned borrowing from Latin \emph{aedicula} 'small shrine', from \emph{aedēs} 'temple' (EtymDB).
\end{cousinweb}

\begin{grammarlens}[title={What you've built}]
One more word for telling, explaining and describing at B1.
\end{grammarlens}

\subsection*{Your turn}
\begin{itemize}
\raggedright
  \item \emph{Say it:} \emph{l'edicolante}
  \item \emph{Say it:} \emph{l'edicolante}, once more
  \item \emph{Say it:} a sentence of your own with \emph{l'edicolante}
  \item \emph{Recall:} say the Italian for the tobacconist, then the Italian for the newsstand, then say \emph{l'edicolante} again
\end{itemize}

\begin{tcolorbox}[breakable,colback=teal!4,colframe=teal!35!black,title={Before you move on}]
What does l'edicolante mean? (\textquotedblleft{}The newsagent\textquotedblright{}.) Say it once more.
\end{tcolorbox}
\section[la cartoleria]{la cartoleria — the stationer's}

\label{lesson:IT-C452-cartoleria}



\begin{quote}
Before the new one: say the Italian for the newsstand, then the Italian for the newsagent.
\end{quote}

\subsection*{You'll want to know: la cartoleria}
\textbf{la cartoleria} — \textquotedblleft{}the stationer's\textquotedblright{}.

\begin{quote}
In cartoleria compro quaderni e penne.
\end{quote}

\begin{quote}
\emph{At the stationer's I buy exercise books and pens.}
\end{quote}

Busy in September, before school starts.

\begin{cousinweb}
\emph{cartoleria} is built on \emph{carta} 'paper', which is inherited from Latin \emph{charta}, from Greek \emph{khártēs} 'papyrus' (EtymDB).
\end{cousinweb}

\begin{grammarlens}[title={What you've built}]
One more word for telling, explaining and describing at B1.
\end{grammarlens}

\subsection*{Your turn}
\begin{itemize}
\raggedright
  \item \emph{Say it:} \emph{la cartoleria}
  \item \emph{Say it:} \emph{la cartoleria}, once more
  \item \emph{Say it:} a sentence of your own with \emph{la cartoleria}
  \item \emph{Recall:} say the Italian for the newsstand, then the Italian for the newsagent, then say \emph{la cartoleria} again
\end{itemize}

\begin{tcolorbox}[breakable,colback=teal!4,colframe=teal!35!black,title={Before you move on}]
What does la cartoleria mean? (\textquotedblleft{}The stationer's\textquotedblright{}.) Say it once more.
\end{tcolorbox}
\section[(dialogue)]{Ripasso dei capitoli 445–452}

\label{lesson:IT-R452-ripasso-445-452}



\begin{quote}
Say the Italian for the bell tower and the main church.
\end{quote}

\subsection*{Your turn}
Say these aloud:

\begin{itemize}
\raggedright
  \item the bell tower, the main church, the cathedral, the arcade, the loggia
  \item the walls, the fortress, the village, the cloister, to wander
  \item the police headquarters, the prefecture, the barracks, to report, the identity card
  \item the bench, the street lamp, the litter bin, the drain cover, to clutter
  \item the block of flats, the scaffolding, the crane, to collapse, to clear
  \item the village food festival, the parade, the procession, the patron saint, the stall
  \item the seafront, to wander around, to browse, to peek, the shop window
  \item the tobacconist's, the tobacconist, the newsstand, the newsagent, the stationer's
\end{itemize}

\begin{tcolorbox}[breakable,colback=teal!4,colframe=teal!35!black,title={Before you move on}]
The bell tower? (\textbf{il campanile}.) The main church? (\textbf{il duomo}.) The cathedral? (\textbf{la cattedrale}.) The arcade? (\textbf{il portico}.) The loggia? (\textbf{la loggia}.) The walls? (\textbf{le mura}.) The fortress? (\textbf{la rocca}.) The village? (\textbf{il borgo}.) The cloister? (\textbf{il chiostro}.) To wander? (\textbf{vagare}.) The police headquarters? (\textbf{la questura}.) The prefecture? (\textbf{la prefettura}.) The barracks? (\textbf{la caserma}.) To report? (\textbf{segnalare}.) The identity card? (\textbf{la carta d'identità}.) The bench? (\textbf{la panchina}.) The street lamp? (\textbf{il lampione}.) The litter bin? (\textbf{il cestino}.) The drain cover? (\textbf{il tombino}.) To clutter? (\textbf{ingombrare}.) The block of flats? (\textbf{il condominio}.) The scaffolding? (\textbf{l'impalcatura}.) The crane? (\textbf{la gru}.) To collapse? (\textbf{crollare}.) To clear? (\textbf{sgomberare}.) The village food festival? (\textbf{la sagra}.) The parade? (\textbf{il corteo}.) The procession? (\textbf{la processione}.) The patron saint? (\textbf{il patrono}.) The stall? (\textbf{la bancarella}.) The seafront? (\textbf{il lungomare}.) To wander around? (\textbf{gironzolare}.) To browse? (\textbf{curiosare}.) To peek? (\textbf{sbirciare}.) The shop window? (\textbf{la vetrina}.) The tobacconist's? (\textbf{la tabaccheria}.) The tobacconist? (\textbf{il tabaccaio}.) The newsstand? (\textbf{l'edicola}.) The newsagent? (\textbf{l'edicolante}.) The stationer's? (\textbf{la cartoleria}.)
\end{tcolorbox}
